% GENERATED FILE. Edit canonical lessons, then run npm run generate:books.
% canonical-source-hash: fnv1a64:bf4ce5df685632f8
% canonical-lessons: MW-C288-mehnat, MW-C288-layak, MW-C288-ustad, MW-C288-dhandho, MW-C288-majur

\chapter{Hard Work}
\label{ch:mw-hard-work}
\hlchaptermodality{288}

\begin{chapteropening}
\textbf{By the end of this chapter:} \emph{I can say what work I do and how hard I work.}

Complete the last lesson of chapter 288: I can say what work I do and how hard I work.
\end{chapteropening}

\section[\texorpdfstring{\mw{मेहनत} (mehnat)}{मेहनत (mehnat)}]{\mw{मेहनत} (mehnat) — hard work}

\label{lesson:MW-C288-mehnat}



\begin{quote}
Before the new one: say the Marwadi for a pupil, then the Marwadi for a wrestling ground.
\end{quote}

\subsection*{You'll want to know: \mw{मेहनत}}
\textbf{\mw{मेहनत}} — \emph{mehnat} — \textquotedblleft{}hard work\textquotedblright{}.

\begin{quote}
\mw{मेहनत} \mw{चोखी} \mw{बात} \mw{है।}
\end{quote}

\begin{quote}
\emph{mehnat chokhī bāt hai.}
\end{quote}

\begin{quote}
Hard work is a good thing.
\end{quote}

A feminine noun.

\begin{cousinweb}
\textbf{\mw{मेहनत}} (\emph{mehnat}) is usually traced, through Persian \emph{mihnat}, \textquotedblleft{}labour\textquotedblright{}, to Arabic \emph{miḥna}, \textquotedblleft{}trial, hardship\textquotedblright{}. \textbf{\mw{मेहनती}}, hard-working, is built on it.
\end{cousinweb}

\begin{grammarlens}[title={What you've built}]
One more word for this chapter.
\end{grammarlens}

\subsection*{Your turn}
\begin{itemize}
\raggedright
  \item \emph{Say it:} \emph{mehnat}
  \item \emph{Say it:} \emph{mehnat}, once more
  \item \emph{Say it:} \emph{mehnat chokhī bāt hai.} — Hard work is a good thing.
  \item \emph{Recall:} say the Marwadi for a pupil, then the Marwadi for a wrestling ground, then say \emph{mehnat} again
\end{itemize}

\begin{tcolorbox}[breakable,colback=teal!4,colframe=teal!35!black,title={Before you move on}]
What does \mw{मेहनत} mean? (\textquotedblleft{}Hard work\textquotedblright{}.) Say it once more.
\end{tcolorbox}
\section[\texorpdfstring{\mw{लायक} (lāyak)}{लायक (lāyak)}]{\mw{लायक} (lāyak) — capable}

\label{lesson:MW-C288-layak}



\begin{quote}
Before the new one: say the Marwadi for a wrestling ground, then the Marwadi for hard work.
\end{quote}

\subsection*{You'll want to know: \mw{लायक}}
\textbf{\mw{लायक}} — \emph{lāyak} — \textquotedblleft{}capable\textquotedblright{}.

\begin{quote}
\mw{छोरो} \mw{घणो} \mw{लायक} \mw{है।}
\end{quote}

\begin{quote}
\emph{chhoro ghaṇo lāyak hai.}
\end{quote}

\begin{quote}
The boy is very capable.
\end{quote}

\begin{cousinweb}
\textbf{\mw{लायक}} (\emph{lāyak}) is usually traced, through Persian, to Arabic \emph{lāyiq}, \textquotedblleft{}worthy, fit\textquotedblright{}.
\end{cousinweb}

\begin{grammarlens}[title={What you've built}]
One more word for this chapter.
\end{grammarlens}

\subsection*{Your turn}
\begin{itemize}
\raggedright
  \item \emph{Say it:} \emph{lāyak}
  \item \emph{Say it:} \emph{lāyak}, once more
  \item \emph{Say it:} \emph{chhoro ghaṇo lāyak hai.} — The boy is very capable.
  \item \emph{Recall:} say the Marwadi for a wrestling ground, then the Marwadi for hard work, then say \emph{lāyak} again
\end{itemize}

\begin{tcolorbox}[breakable,colback=teal!4,colframe=teal!35!black,title={Before you move on}]
What does \mw{लायक} mean? (\textquotedblleft{}Capable\textquotedblright{}.) Say it once more.
\end{tcolorbox}
\section[\texorpdfstring{\mw{उस्ताद} (ustād)}{उस्ताद (ustād)}]{\mw{उस्ताद} (ustād) — a master of a craft}

\label{lesson:MW-C288-ustad}



\begin{quote}
Before the new one: say the Marwadi for hard work, then the Marwadi for capable.
\end{quote}

\subsection*{You'll want to know: \mw{उस्ताद}}
\textbf{\mw{उस्ताद}} — \emph{ustād} — \textquotedblleft{}a master of a craft\textquotedblright{}.

\begin{quote}
\mw{म्हारो} \mw{उस्ताद} \mw{घणो} \mw{हुसियार} \mw{है।}
\end{quote}

\begin{quote}
\emph{mhāro ustād ghaṇo husiyār hai.}
\end{quote}

\begin{quote}
My master is very clever.
\end{quote}

A masculine noun.

\begin{cousinweb}
\textbf{\mw{उस्ताद}} (\emph{ustād}) is usually traced to Persian \emph{ustād}, \textquotedblleft{}a master\textquotedblright{}: the one who has the skill and teaches it.
\end{cousinweb}

\begin{grammarlens}[title={What you've built}]
One more word for this chapter.
\end{grammarlens}

\subsection*{Your turn}
\begin{itemize}
\raggedright
  \item \emph{Say it:} \emph{ustād}
  \item \emph{Say it:} \emph{ustād}, once more
  \item \emph{Say it:} \emph{mhāro ustād ghaṇo husiyār hai.} — My master is very clever.
  \item \emph{Recall:} say the Marwadi for hard work, then the Marwadi for capable, then say \emph{ustād} again
\end{itemize}

\begin{tcolorbox}[breakable,colback=teal!4,colframe=teal!35!black,title={Before you move on}]
What does \mw{उस्ताद} mean? (\textquotedblleft{}A master of a craft\textquotedblright{}.) Say it once more.
\end{tcolorbox}
\section[\texorpdfstring{\mw{धंधो} (dhandho)}{धंधो (dhandho)}]{\mw{धंधो} (dhandho) — a trade}

\label{lesson:MW-C288-dhandho}



\begin{quote}
Before the new one: say the Marwadi for capable, then the Marwadi for a master of a craft.
\end{quote}

\subsection*{You'll want to know: \mw{धंधो}}
\textbf{\mw{धंधो}} — \emph{dhandho} — \textquotedblleft{}a trade\textquotedblright{}.

\begin{quote}
\mw{म्हारो} \mw{धंधो} \mw{सोरो} \mw{कोनी।}
\end{quote}

\begin{quote}
\emph{mhāro dhandho soro konī.}
\end{quote}

\begin{quote}
My trade isn't easy.
\end{quote}

A masculine noun.

\begin{cousinweb}
\textbf{\mw{धंधो}} (\emph{dhandho}) is inherited from Sanskrit \emph{dhandha}, \textquotedblleft{}awkwardness\textquotedblright{}; CDIAL gives no Marwari form in this entry, but lists Hindi \emph{dhandhā}, \textquotedblleft{}craft, way of life\textquotedblright{}, and Sindhi \emph{dhandho}, \textquotedblleft{}business\textquotedblright{} (CDIAL 6727). How the sense moved to work is not settled.
\end{cousinweb}

\begin{grammarlens}[title={What you've built}]
One more word for this chapter.
\end{grammarlens}

\subsection*{Your turn}
\begin{itemize}
\raggedright
  \item \emph{Say it:} \emph{dhandho}
  \item \emph{Say it:} \emph{dhandho}, once more
  \item \emph{Say it:} \emph{mhāro dhandho soro konī.} — My trade isn't easy.
  \item \emph{Recall:} say the Marwadi for capable, then the Marwadi for a master of a craft, then say \emph{dhandho} again
\end{itemize}

\begin{tcolorbox}[breakable,colback=teal!4,colframe=teal!35!black,title={Before you move on}]
What does \mw{धंधो} mean? (\textquotedblleft{}A trade\textquotedblright{}.) Say it once more.
\end{tcolorbox}
\section[\texorpdfstring{\mw{मजूर} (majūr)}{मजूर (majūr)}]{\mw{मजूर} (majūr) — a labourer}

\label{lesson:MW-C288-majur}



\begin{quote}
Before the new one: say the Marwadi for a master of a craft, then the Marwadi for a trade.
\end{quote}

\subsection*{You'll want to know: \mw{मजूर}}
\textbf{\mw{मजूर}} — \emph{majūr} — \textquotedblleft{}a labourer\textquotedblright{}.

\begin{quote}
\mw{मजूर} \mw{खेत} \mw{मांय} \mw{मेहनत} \mw{करै} \mw{है।}
\end{quote}

\begin{quote}
\emph{majūr khet mā̃y mehnat karai hai.}
\end{quote}

\begin{quote}
The labourer works hard in the field.
\end{quote}

A masculine noun.

\begin{cousinweb}
\textbf{\mw{मजूर}} (\emph{majūr}) is usually traced to Persian \emph{mazdūr}, \textquotedblleft{}a labourer\textquotedblright{}.
\end{cousinweb}

\begin{grammarlens}[title={What you've built}]
One more word for this chapter.
\end{grammarlens}

\subsection*{Your turn}
\begin{itemize}
\raggedright
  \item \emph{Say it:} \emph{majūr}
  \item \emph{Say it:} \emph{majūr}, once more
  \item \emph{Say it:} \emph{majūr khet mā̃y mehnat karai hai.} — The labourer works hard in the field.
  \item \emph{Recall:} say the Marwadi for a master of a craft, then the Marwadi for a trade, then say \emph{majūr} again
\end{itemize}

\begin{tcolorbox}[breakable,colback=teal!4,colframe=teal!35!black,title={Before you move on}]
What does \mw{मजूर} mean? (\textquotedblleft{}A labourer\textquotedblright{}.) Say it once more.
\end{tcolorbox}
